\backhead{FURTHER EXERCISES}{khred}
\begin{exlist}
\item For each of the ten soils, is a 600 cm trigger inside the reach? List the soils for which it never fires.\exans{sand and loamy sand}
\item Using the table of readily available water, rank the soils by the interval between irrigations at $ET_c=5$ mm\,d$^{-1}$.\exans{sand shortest, silt loam longest}
\item Find the pore radius that is just full at pF 3.0 and at pF 4.2.\exans{1.5 $\mu$m; 0.094 $\mu$m}
\item For the loam, find $S_e$ and $K_r$ at 1000 cm.\exans{0.1342; $6.55\times10^{-7}$}
\item Convert 300 cm of suction to kPa, bar and pF.\exans{29.4 kPa; 0.294 bar; 2.48}
\end{exlist}
\backhead{ESSAY QUESTIONS}{khred}
{\small\itshape Answer each question in about one page, with equations, a sketch and at least one number from this chapter.}\par\vspace{3pt}
\begin{exlist}
\item Choose two soils of contrasting texture and compare their worked solutions step by step. What do the differences mean for irrigation?
\item Explain how the reference tables of this chapter can be used to design an irrigation schedule without a computer.
\item Discuss the accuracy of the reference tables and when a full computation is needed instead.
\item Explain the link between the pore radius, the suction and the pF.
\item Discuss the reach of the trigger for the ten soils and what it means for advice to farmers.
\item Write a one-page guide to irrigating the loam, using only the tables of this chapter.
\end{exlist}
